%! Author = drakanoy
%! Date = 10.09.2024

% Preamble
\documentclass[12pt]{article}

% Packages
\usepackage[utf8]{inputenc}
\usepackage[T2A]{fontenc}
\usepackage[english, russian]{babel}
\usepackage[a4paper, includefoot, left=1.5cm, right=1.5cm, top=1cm, bottom=1.5cm, headsep=1cm, footskip=1cm]{geometry}
\usepackage{makecell}
\usepackage{amsmath}
\usepackage{graphicx}
\usepackage{enumitem}
\usepackage{svg}
\usepackage{multirow}
\usepackage{hyperref}
\usepackage{mathtools}
\usepackage{amssymb}
\usepackage{textcomp}
\usepackage{stmaryrd}

% Document
\begin{document}
\begin{large}
\begin{center}
\LARGE \textbf{Оптика Полупроводников}
\par
\LARGE \textbf{Кононов Александр Михайлович}
\par
    \textbf{18.02.2026}
\end{center}
\par \textbf{Задача}:
\par Условие:
\par В среде с диэлектрической проницаемостью $\varepsilon$ (с показателем приломления $n$) колеблится диполь $d$ с частотой $\omega_0$.
Найти темп убыли энергии и сравнить с результатом полученным на лекции.
\par
\includegraphics[width=0.5\textwidth]{photo_1.jpg}
\begin{center}
\underline{Рисунок 1}:
\end{center}
\par \textbf{Решение}:
%\par
%\includegraphics[width=1\textwidth]{photo_11.jpg}
%\par
\par Из запаздывающего потенциала в волновой области пространства верны следующие соотношения на векторный потенциал и магнитное поле:
\[
    \vec{A} = \frac{\dot{\vec{d}}}{cR}
\]
\[
    \vec{H} = rot \vec{A} = \frac{1}{c} \left[ \dot{\vec{A}} \times \vec{n} \right]
\]
где $\vec{n}$ - единичный вектор направления.
\par Подставляем и получаем:
\[
    \vec{A} = \frac{\omega_0 \cdot \cos\left( \omega_0 t \right)}{cR} \vec{d_0}
\]
\[
    \dot{\vec{A}} = -\frac{d_0 \omega_0^2 \cdot \sin\left( \omega_0 t \right)}{cR} \vec{e_z}
\]
\par Посчитаем магнитное поле:
\[
    \vec{H} = \frac{d_0 \omega_0^2 \cdot \sin\left( \omega_0 t \right)}{c^2 R} \left[ \vec{n} \times \vec{e_z} \right] = -\frac{d_0 \omega_0^2 \cdot \sin\left( \omega_0 t \right)}{c^2 R} \sin\theta \cdot \vec{e_{\phi}}
\]
\par Интенсивность излучения в единице телесного угла:
\[
    dI = \left< \left| \vec{s} \right| \right>_t R^2 d\Omega = \frac{nc}{4\pi} \left< \vec{H}^2 \right>_t R^2 d\Omega
\]
\par Подставим магнитное поле сразу усреднив $\sin^2$ до $1/2$:
\[
    dI = \frac{nc}{4\pi} \frac{\omega_0^4 d_0^2}{2c^4} \sin^2\theta d\Omega = \frac{n\omega_0^4 d_0^2}{8\pi c^3} \sin^2\theta d\Omega
\]
\par Проинтегрировав по телесному углу, получим:
\[
    I = \frac{n\omega_0^4 d_0^2}{8\pi c^3} \cdot 2\pi \frac{4}{3} = \frac{n\omega_0^4 d_0^2}{3c^3}
\]
\par Интенсивность излучения это убыль энергии в единицу времени. Считая что вся энергия высветилась в один фотон с частотой $\omega_0$ за время $\tau$, получим темп перехода:
\[
    \frac{1}{\tau} = \frac{I}{E} = \frac{1}{\hslash \omega_0} \cdot \frac{n\omega_0^4 d_0^2}{3c^3} = \frac{n\omega_0^3 d_0^2}{3 \hslash c^3}
\]
\par Получился тот же ответ, что и на лекциях с точностью до коэффициента 4.
Это можно объяснить так: на лекциях мы для вектора Пойнтинга использовали формулу $<S> = \frac{nc}{2\pi}E^2$, когда как при решении задачи я использовал $<S> = \frac{nc}{4\pi}E^2$ - в 2 раза меньше.
Вторую двойку можно объяснить тем, что по закону сохранения импульса диполь должен высветить 2 фотона в противоположные стороныБ что увеличит интенсивность в 2 раза и мы получим тот же ответ, что и на лекциях.
%\par \textbf{Ответ}:
%\par Возраст Земли, рассчитанный на основе долей разных изотопов урана, составил $T = 5,94 \cdot 10^9$ лет.
%\includegraphics[width=1\textwidth]{photo_1.jpg}
%\par
%\begin{center}
%\includegraphics[width=0.4\textwidth]{photo_1.jpg}
%\end{center}

\newpage

\[
    \hat{H} = g \left( \hat{b}^\dag_{-} \hat{a}_{+} + \hat{b}_{-} \hat{a}^\dag_{+} \right)
\]
\[
    \left| \psi_{in} \right> = e^{-\left| \alpha \right|^2 / 2} e^{\alpha \hat{b}^{\dag}_H} \left| 0 \right> = e^{-\left| \alpha \right|^2 / 2} \sum_n \frac{\alpha^n}{\sqrt{n!}}  \left| n_H \right>
\]
\[
    \left| \psi_{in}^H \right> =  e^{\alpha \hat{b}^{\dag}_H} \left| 0 \right> = \sum_n \frac{\alpha^n}{\sqrt{n!}}  \left| n_H \right>
\]


\begin{gather*} \nonumber
    \left| \psi_{out} \right> = \frac{e^{-\left| \alpha \right|^2 / 2} \ e^{\alpha \hat{b}^{\dag}_{H}}}{2}  \left[ e^{-i \theta / 2 } \ e^{-i \gamma \hat{b}^{\dag}_{-} } + e^{i \theta / 2 } \ e^{i \gamma \hat{b}^{\dag}_{-} } + \\
    + \left. \left( e^{-i \theta / 2 } \ e^{-i \gamma \hat{b}^{\dag}_{-} } - e^{i \theta / 2 } \ e^{i \gamma \hat{b}^{\dag}_{-} } \right) \hat{a}^{\dag}_{+} \right] \left| 0 \right>
\end{gather*}

\[
    \left| \psi_{out}^V \right> \approx \cos \left( \frac{\theta }{2} \right) \left| 0 \right> + \frac{i \delta - 1}{\sqrt{2}} \sin \left( \frac{\theta}{2} \right) \left| 1 \right> + \frac{i \delta}{2} \cos \left( \frac{\theta}{2} \right) \left| 2 \right>
\]
\[
    \left| \psi_{out} \right> = \hat{S_0} \left| \psi_{in} \right> =  e^{\alpha \hat{b}^{\dag}_{+} / \sqrt{2}} \ \hat{S_0}  \ e^{\alpha \hat{b}^{\dag}_{-} / \sqrt{2}} \ \left| 0 \right>
\]
\[
    i \frac{\partial \hat{\rho}_{HV} }{\partial t} = \left[ \hat{H} ; \hat{\rho}_{HV} \right]
\]
\[
     \hat{\rho}_{HV} = e^{\alpha \hat{b}^{\dag}_{H}}  \hat{\rho}_{V}  e^{\alpha^{*} \hat{b}_{H}}
\]
\[
    i \frac{\partial \hat{\rho}_{V} }{\partial t} = \left[ \hat{H}_V ; \hat{\rho}_{V} \right]
\]
\[
    \hat{H}_V = \Omega_R^{H} \left( \hat{d}^{\dag}_{\uparrow} \hat{a}_{\uparrow}
    + \hat{d}_{\uparrow} \hat{a}^{\dag}_{\uparrow} \right)
    + ig \left( \hat{b}_{V} \hat{d}^{\dag}_{\uparrow} \hat{a}_{\uparrow}
    - \hat{b}^{\dag}_{V} \hat{d}_{\uparrow} \hat{a}^{\dag}_{\uparrow} \right)
\]
\[
    g << \Omega_R^{H}
\]

\par
\end{large}
\end{document}
